\documentclass{article}
\usepackage{amsmath, amssymb, amsfonts}
\title{X-ray Diffraction and Reciprocal Lattice}
\begin{document}
\abstract{David Snoke introduces the reciprocal lattice, the structure factor
and their relation with X-ray diffraction far more cleanly than Charles Kittel.
This article is a summary of his treatment.}
\author{Amey Joshi}
\maketitle
\section{The reciprocal lattice}
Consider a crystal with translation vectors
\begin{equation}\label{e1}
\vec{R} = u_1\vec{a}_1 + u_2\vec{a}_2 + u_3\vec{a}_3,
\end{equation}
where $u_1, u_2, u_3 \in \mathbb{Z}$ and $\vec{a}_i$ are the primitive lattice 
vectors. Suppose that at each lattice point, there is a basis consisting of $s$
atoms whose positions relative to it are given by $\vec{r}_i$. Then the crystal
can be described by the function
\begin{equation}\label{e2}
f(\vec{r}) = \sum_{\vec{R}}\sum_{i=1}^s \delta(\vec{r} - \vec{R} - \vec{r}_i).
\end{equation}
Consider the Fourier transform of $f$,
\begin{equation}\label{e3}
F(\vec{k}) = \int_{-\infty}^\infty f(\vec{r})e^{i\vec{k}\cdot\vec{r}}d\vec{r}
= \sum_{\vec{R}}\sum_{i=1}^se^{i\vec{k}\cdot(\vec{R} + \vec{r}_i)}.
\end{equation}
We can separate the basis terms and write $F$ as
\begin{equation}\label{e4}
F(\vec{k}) = \sum_{\vec{R}}e^{i\vec{k}\cdot\vec{R}}\sigma,
\end{equation}
where
\begin{equation}\label{e5}
\sigma = \sum_{i=1}^s e^{i\vec{k}\cdot\vec{r}_i},
\end{equation}
is the structure factor of the crystal.
The function $F$ vanishes for all $\vec{k}$ except for those that satisfy the
condition
\begin{equation}\label{e6}
\exp(i\vec{k}\cdot\vec{R}) = 0.
\end{equation}
This can happen only when $\vec{k}\cdot\vec{R} = 2\pi N$, for some $N \in
\mathbb{Z}$. A $\vec{k}$ that satisfies this condition is of the form
\begin{equation}\label{e7}
\vec{K} = v_1\vec{b}_1 + v_2\vec{b}_2 + v_3\vec{b}_3,
\end{equation}
where $v_1, v_2, v_3 \in \mathbb{Z}$ and 
\begin{eqnarray}
\vec{b}_1 &=& 2\pi\frac{\vec{a_2}\times\vec{a}_3}
{\vec{a}_1\cdot(\vec{a}_2\times\vec{a}_3)} \label{e8} \\
\vec{b}_2 &=& 2\pi\frac{\vec{a_3}\times\vec{a}_1}
{\vec{a}_2\cdot(\vec{a}_3\times\vec{a}_1)} \label{e9} \\
\vec{b}_3 &=& 2\pi\frac{\vec{a_1}\times\vec{a}_2}
{\vec{a}_3\cdot(\vec{a}_1\times\vec{a}_2)} \label{e10} 
\end{eqnarray}
The vectors $\vec{b}_i$ form the primitive vectors of the reciprocal lattice.
The vector $\vec{K}$ is then a translation vector of the reciprocal lattice.


\section{X-ray diffraction}
We consider Fraunhoffer diffraction of X-rays due to atoms in a crystal. That
is, the incident and diffracted waves are plane waves and are measured far from
the diffraction site. If the incident wave is
\begin{equation}\label{e11}
A_i(\vec{r}) = A_0e^{i(\vec{k}_0\cdot\vec{r} - \omega t)},
\end{equation}
then the diffracted wave is of the form
\begin{equation}\label{e12}
A_d(\vec{r}) = A_0e^{i(\vec{k}\cdot\vec{r} - \omega t)},
\end{equation}
where the diffraction only results in scattering the wave by an amount
\begin{equation}\label{e13}
\vec{s} = \vec{k} - \vec{k}_0.
\end{equation}
As a result, we can write \eqref{e12} as
\begin{equation}\label{e14}
A_d(\vec{r}) = A_i e^{i\vec{s}\cdot\vec{r}}.
\end{equation}
The wave diffracted from the atoms at $\vec{r} + \vec{R}$ will be
\begin{equation}\label{e15}
A_d(\vec{r} + \vec{R}) = A_d(\vec{r})e^{i\vec{s}\cdot\vec{R}}.
\end{equation}
Waves diffracted by all points in the crystal are superposed to form the 
resultant wave
\begin{equation}\label{e16}
A_d = \sum_{\vec{R}} A_d(\vec{r} + \vec{R}) = A_d(\vec{r})
\sum_{\vec{R}} e^{i\vec{s}\cdot\vec{R}}.
\end{equation}
$A_d$ will generally be zero except for those $\vec{s}$ for which the summation
over all lattice points does not vanish. In the previous section, we concluded
that this can happen only if $\vec{s}$ are translation vectors of the reciprocal
lattice. Thus, the condition for diffraction is
\begin{equation}\label{e17}
\Delta\vec{k} = \vec{k} - \vec{k}_0 = \vec{s} = \vec{G},
\end{equation}
where $\vec{G}$ is a translation vector in the reciprocal lattice. This is
equation (21) in Kittel's second chapter.

In general, it is hard to satisfy \eqref{e17} for all wavelengths of the X-rays
and all directions of scattering. There are two ways to achive the diffraction
condition in practice.
\subsection{Laue diffraction}
In this case the source of X-rays has a broad spectrum and the direction of the
beam is held fixed. Align the $z$-axis with the incident beam. Therefore, the
incident wave vector is $\vec{k}_0 = k_0\hat{z}$. We assume that the scattering
is elastic. Therefore, the scattered vector will have the same magnitude as the
incident one so that
\begin{equation}\label{e18}
\vec{k} = k_0(\sin\theta\cos\phi\hat{x} + \sin\theta\sin\phi\vec{y} + 
\cos\theta\hat{z})
\end{equation}
and
\begin{equation}\label{e19}
\vec{s} = k_0(\sin\theta\cos\phi\hat{x} + \sin\theta\sin\phi\vec{y} + 
(\cos\theta - 1)\hat{z})
\end{equation}
For a given orientation of the detector, the two angles are fixed, but the
broad spectrum of the incident radiation results in a number of spots, each
one for the wavelength that satisfies equation \eqref{e19}.

\subsection{Powder diffraction}
In this method, the solid is crushed into a fine powder and the incident beam is
monochromatic. If $\lambda = 2\pi/k_0$ is the wavelength of the incident beam
then we are guaranteed that the diffracted beam has the same wavelength. If $
\alpha$ is the angle between $\vec{k}_0$ and $\vec{k}$ then the diffraction
condition $\vec{s} = \vec{G}$ gives $s^2 = G^2$ or $2k_0^2(1 + \cos\alpha) = 
G^2$. That is
\[
4k_0^2\sin^2\frac{\alpha}{2} = G^2
\]
or
\begin{equation}\label{e20}
\frac{4\pi}{\lambda}\sin\frac{\alpha}{2} = |\vec{G}|.
\end{equation}
Each $\vec{G}$ results in a bright ring caused by X-rays diffracted in a cone
of semi-angle $\alpha$ about the incident beam. The cone intersects the film in
a circle.

\end{document}
